% title:          Commutative subgroups generated by powers
% area:           group-theory, free-groups
% methods:        induction, bezout
% difficulty:
% difficulty-ai:  hard
% status:
% review:         none
% --- history: all optional, leave EMPTY if unknown ---
% origin:         imc
% origin-date:    2005-07
% origin-number:  Day 1, Problem 6
% also-in:
% taken-from:
% references:     IMC 2005 (12th IMC, Blagoevgrad, 22-28 July 2005), First Day, Problem 6, official problems & solutions - https://www.imc-math.org.uk/imc2005/day1_solutions.pdf; IMC 2005 problems index page - https://www.imc-math.org.uk/?year=2005&item=problems; later reprinted as Problem 2.17 '(2005)' in M. A. Daas, 'A Brief IMC Training' (Mathematical Institute Leiden, 2024) - http://www.cs.emory.edu/~mic/putnam/general/IMC.pdf
% history:        ai
\begin{problem}
Let $m,n\in\mathbb{Z}$. Given a group \( G \), denote by \( G(m) \) the subgroup generated by the \( m \)-th powers of elements of \( G \):
\[
G(m):=\left\langle\left\{g^m \mid g\in G\right\}\right\rangle\leq G.
\]
If \( G(m) \) and \( G(n) \) are commutative, prove that \( G(\gcd(m, n)) \) is also commutative. Here, \( \gcd(m, n) \) denotes the greatest common divisor of \( m \) and \( n \).
\end{problem}
\begin{solution}
If $m=0$ or $n=0$, this is trivial. Suppose now $|m|,|n|\geq 1$.
\\\\
Recall that if $H$ is a group and $H'\subset H$ is a subset, then the subgroup of $H$ generated by $H'$ is the smallest subgroup of $H$ containing $H'$, that is, $\langle H'\rangle=\bigcap_{H'\subset F\leq H}F$. It is easy to see that:
\begin{equation}\label{1}
\langle H'\rangle=\left\{\prod_{i\in k}g_i^{\epsilon_i} \,\Bigg|\, \exists k\in\mathbb{N}\, \exists g\in (H')^k\, \exists \epsilon\in\{-1,1\}^{k}\right\},
\end{equation}
where the product is taken in the order of the integer $k=\{i\in\mathbb{N} \mid i<k\}$ and the empty product is $e_H$.
\\\\
Write \( d = \gcd(m, n) \). Notice that:
\[
G(d)=\left\langle G(m)\cup G(n) \right\rangle.
\]
Indeed, this follows from the monotonicity of $\langle\_\rangle$ and the fact that every subgroup is a fixed point of $\langle\_\rangle$. If $z\in\left\{g^d \mid g\in G\right\}$, then $z=g^d$ for some $g\in G$. By Bézout's lemma, there exist two integers $l,r\in\mathbb{Z}$ with $lm+rn=d$, and we have:
\[
z=g^d=g^{lm+rn}=\left(\left(g^m\right)^{\mathrm{sign}(l)}\right)^{|l|}\left(\left(g^n\right)^{\mathrm{sign}(r)}\right)^{|r|}.
\]
Because $g^m\in G(m)$ and $g^n\in G(n)$, we have $z\in\left\langle G(m)\cup G(n) \right\rangle$. Since $z$ was arbitrary, we get $\left\{g^d \mid g\in G\right\}\subset\left\langle G(m)\cup G(n)\right\rangle$, and thus:
\[
G(d)=\left\langle\left\{g^d \mid g\in G\right\}\right\rangle\subset \left\langle G(m)\cup G(n)\right\rangle.
\]
Similarly, let $z\in \left\{g^m \mid g\in G\right\}$. Then $z=g^m$ for some $g\in G$, and thus:
\[
z=\left(g^{d}\right)^{\frac{m}{d}}\in G(d),
\]
since $g^d\in G(d)$. As $z$ was arbitrary, we conclude $\left\{g^m \mid g\in G\right\}\subset G(d)$. Thus, $G(m)\subset G(d)$. In the exact same manner, $G(n)\subset G(d)$, and therefore:
\[
\left\langle G(m)\cup G(n)\right\rangle\subset G(d).
\]
We also have:
\[
\left\langle G(m)\cup G(n) \right\rangle=\left\langle\left\{g^m \mid g\in G\right\}\cup\left\{h^n \mid h\in G\right\}\right\rangle\footnote{If $H$ is a group, $I$ a set, and $\{H_i \mid i\in I\}$ are subgroups of $H$ each generated by $\{S_i \mid i\in I\}\subset\mathscr{P}(H)$ respectively ($\forall i\in I, H_i=\langle S_i\rangle$), then $\left\langle\bigcup_{i\in I} H_i\right\rangle=\left\langle\bigcup_{i\in I}S_i\right\rangle$. Indeed, we clearly have $\bigcup_{i\in I}S_i\subset\bigcup_{i\in I}H_i$ (since $S_i\subset \langle S_i\rangle$) and thus $\left\langle\bigcup_{i\in I}S_i\right\rangle\subset \left\langle\bigcup_{i\in I} H_i\right\rangle$. If we take an element $z\in \bigcup_{i\in I}H_i$, then there exists $j\in I$ with $z\in H_j=\langle S_j\rangle\subset\left\langle\bigcup_{i\in I}S_i\right\rangle$, and so we have $\bigcup_{i\in I}H_i\subset\left\langle\bigcup_{i\in I}S_i\right\rangle$, which means that we have the other inclusion $\left\langle\bigcup_{i\in I}H_i\right\rangle\subset\left\langle\bigcup_{i\in I}S_i\right\rangle$.}.
\]
It is also clear regarding equation (\ref{1}) that if $S\subset G$ is constituted of elements that commute with one another, then $\langle S\rangle$ is a commutative subgroup (this follows from the fact that if $a,b\in S$ commute, then any two elements in $\{a^{-1},b^{-1},a,b\}$ commute). The converse of this statement is trivial. Therefore, by the two equalities above, showing commutativity of $G(d)$ is equivalent to showing commutativity of any two elements in $\left\{g^m \mid g\in G\right\}\cup\left\{h^n \mid h\in G\right\}$. Because we know that any two elements in $\left\{g^m \mid g\in G\right\}$ or $\left\{h^n \mid h\in G\right\}$ commute (since $G(m)$ and $G(n)$ are commutative), we only need to show that any element in $\left\{g^m \mid g\in G\right\}$ commutes with any other element in $\left\{h^n \mid h\in G\right\}$. So, without further ado, let any two generators \( a^m \) and \( b^n \) ($a,b\in G$). Showing commutativity of these two elements is equivalent to showing neutrality of their commutator:
\[
z:= [a^m,b^n] = a^{-m}b^{-n} a^m b^n.
\]
Then the relations
\[
z = \left(a^{-m} b a^m\right)^{-n} b^n = a^{-m} \left(b^{-n} a b^n\right)^m
\]
show that \( z \in G(m) \cap G(n) \). But then \( z \) is in the center of \( G(d) \). Indeed to show $z\in Z(G(d))$, we let $x\in G(d)=\left\langle\left\{g^m \mid g\in G\right\}\cup\left\{h^n \mid h\in G\right\}\right\rangle$, then there exists a natural number $k$ and a sequence of length $2k$ of element of the group $\textbf{g}\in G^{2k}$ with:
\[
x=g_0^m g_1^n \cdots g_{2k-2}^m g_{2k-1}^n,
\]
and as $z\in G(m) \cap G(n)$, $z$ commutes with any element in $\left\{g^m \mid g\in G\right\}\cup\left\{h^n \mid h\in G\right\}$ so we have by induction $zx=xz$. Now, from the relation \( a^m b^n = b^n a^m z \), it easily follows by induction that for any integer $l\geq 0$, we have:
\[
a^{ml} b^{nl}=b^{nl} a^{ml} z^{l^2}.
\]
Indeed, for $l=0$, we have:
\[
a^{ml} b^{nl}=a^{0} b^{0}=e_G=b^{0} a^{0} z^0=b^{nl} a^{ml} z^{l^2},
\]
and for $l=1$, we have:
\[
a^{ml} b^{nl}=a^{m} b^{n}=b^n a^m z=b^{nl} a^{ml} z^{l^2},
\]
where we used the above relation. Suppose this holds for $l\geq 1$. We show this holds for $l+1$. We have:
\[
a^{m(l+1)} b^{n(l+1)}=a^m \left(a^{ml}b^{nl}\right)b^n=a^m b^{nl} a^{ml} z^{l^2} b^n=a^m b^{nl} a^{ml} b^n z^{l^2},
\]
where we use in the second equality the induction hypothesis and in the third the fact that $z$ is in the center and hence $z^{l^2}$ as well. Now:
\[
a^m b^{nl}=b^n a^m z b^{n(l-1)}=b^n \left(a^m b^{n(l-1)}\right) z=\overset{l-1\text{ times}}{\cdots}=b^{nl} a^{m} z^{l},
\]
where we do an induction on $l$ by iteratively using the known relation and the fact that $z$ is in the center. Thus:
\[
a^{m(l+1)} b^{n(l+1)}=a^m b^{nl} a^{ml} b^n z^{l^2}=b^{nl} a^{m(l+1)} b^n z^{l^2+l},
\]
where we used the two last results and the fact that $z^l$ is in the center. Similarly:
\[
a^{m(l+1)} b^{n}=\left(a^{ml} b^n\right) a^m z=\overset{l\text{ times}}{\cdots}=b^{n} a^{m(l+1)} z^{l+1},
\]
where we do an induction on $l+1$ by iteratively using the known relation and the fact that $z$ is in the center. Thus:
\[
a^{m(l+1)} b^{n(l+1)}=b^{nl} a^{m(l+1)} b^n z^{l^2+l}=b^{n(l+1)} a^{m(l+1)} z^{l^2+2l+1}=b^{n(l+1)} a^{m(l+1)} z^{(l+1)^2},
\]
where we used the two last results and $l^2+2l+1=(l+1)^2$. This concludes the induction and proves the statement.
\\\\
In particular, for any integer $l\geq 0$, we have:
\[
z^{l^2}=a^{-ml} b^{-nl} a^{ml} b^{nl}=[a^{ml}, b^{nl}].
\]
Setting the two integers \( k := \frac{m}{d} \) and $k':=\frac{n}{d}$, since $nk=mk'$, we obtain that $a^{mk}=(a^k)^{m}$ and $b^{nk}=(b^{k'})^{m}$ i.e. $a^{mk},b^{nk}\in G(m)$, and thus they commute by hypothesis, which means:
\[
z^{k^2}=[a^{mk}, b^{nk}] = e_G.
\]
Similarly, $a^{mk'}=(a^k)^{n}$ and $b^{nk'}=(b^{k'})^{n}$ i.e. $a^{mk'},b^{nk'}\in G(n)$, and thus they commute by hypothesis, which means:
\[
z^{k'^2}=[a^{mk'}, b^{nk'}] = e_G.
\]
So $z^{\left(\frac{m}{d}\right)^2}=e_G=z^{\left(\frac{n}{d}\right)^2}$. Clearly, $\gcd(k,k')=1$, and thus $\gcd(k^2,k'^2)=1$, so by Bézout's lemma, there exist two integers $s,t\in\mathbb{Z}$ with $s k^2 + t k'^2=1$. Hence:
\[
e_G=(e_G)^s (e_G)^t=\left(z^{k^2}\right)^s \left(z^{k'^2}\right)^t=z^{s k^2 + t k'^2}=z^1=[a^m, b^n].
\]
Since $a,b\in G$ were arbitrary, we conclude that \( G(d) \) is commutative, as required.
\end{solution}
